\section{Evaluation}
\label{sec:evaluation}

\textbf{Setup.}
We use Qwen3-8B in BF16 on two Ascend 910B3 NPUs (64\,GB each) with TP=2,
CANN 8.3.rc1, PyTorch/torch-npu 2.7.1, and vLLM-Ascend v0.11.0rc1. The direct
baseline, \baseline{}, and PD-TDM use the same replica and 2,048-token prefill
bound. T6 replays Azure-derived conversation and code bursts with a 10-s cycle.
Mean prompt/output lengths are 1,155/211 and 2,048/28 tokens, respectively.
Each point runs for 90\,s with a 30\,s warm-up; we report the median of three
arrival seeds. Strict TTFT/TPOT thresholds are 200/120\,ms for conversation and
500/200\,ms for code. Metrics include requests whose arrivals fall in the
measurement cohort. We use only runs collected after fixing the scheduler's
waiting gate.

\begin{figure*}[t]
  \centering
  \includegraphics[width=0.91\textwidth]{fig_t6_strict_slo.pdf}
  \caption{Strict joint-SLO attainment and output-token goodput on
  Azure-derived bursts. Each point is the median of three seeds. vLLM-Ascend CP
  and PD-TDM use the same two-NPU TP replica and 2,048-token prefill bound.}
  \label{fig:t6-main}
\end{figure*}

\textbf{Azure bursts.}
Fig.~\ref{fig:t6-main} shows the load sweep. At conversation scale $k=3.0$,
PD-TDM raises attainment from 29.0\% to 63.8\% (+34.8 points) and goodput from
1,525 to 3,396 output token/s. At code scale $k=4.9$, attainment rises from
6.5\% to 80.4\% (+73.9 points) and goodput from 87 to 438 output token/s.
The large relative goodput at this second point should be read with its low
baseline attainment, not as a workload-independent speedup. Wider SLOs often
place both systems near 100\% attainment, indicating that the main effect is a
latency-distribution shift rather than unconditional raw-throughput growth.

Scheduler telemetry supports the execution claim without establishing a kernel
cause. In one persistently backlogged code point, PD-TDM alternates pure $P/D$
iterations and the decode-iteration inter-arrival P50/P99 is 189/203\,ms. Mixed
and pure iterations perform unequal work, so their durations do not establish a
fixed mixed-kernel tax.

\begin{table}[t]
\centering
\caption{Supporting MaaS replay results; latency in ms.}
\label{tab:maas}
\small
\begin{tabular}{lrr}
\toprule
Metric & \baseline{} & PD-TDM \\
\midrule
TTFT mean / P99 & 952 / 1,316 & 808 / 1,180 \\
TPOT mean / P99 & 236 / 257 & 222 / 252 \\
$3\times$ SLO attainment & 99.997\% & 99.997\% \\
$1.2\times$ attainment$^{*}$ & 63.4\% & 92.6\% \\
$1.2\times$ goodput$^{*}$ & 99.9 & 150.5 \\
\bottomrule
\multicolumn{3}{l}{\footnotesize $^{*}$Post-hoc; goodput in output token/s.}
\end{tabular}
\end{table}

\textbf{MaaS replay and boundary.}
We reconstruct request lengths from minute-level aggregates and replay a
76-minute arrival trace containing 35,441 submitted requests per system, with
no client errors. The configured cohort window is $[30,34200)$\,s; it contains
35,345 arrivals, of which 21 baseline and 3 PD-TDM records lack matched latency
telemetry. Table~\ref{tab:maas} reports matched-record latency summaries.
PD-TDM lowers mean/P99 TTFT by 15.1\%/10.4\%. Under the original loose
$3\times$ ideal SLO, both systems are at the ceiling. Tightening the same traces
to $1.2\times$ ideal produces the listed separation, but this is a post-hoc
sensitivity analysis rather than a prespecified service target.

Synthetic workloads expose the other side of the tradeoff. Balanced loads show
TTFT gains, whereas longer-output, decode-heavy cells become neutral and then
favor the mixed baseline at high load. These boundary runs are single seed and
some high-load cells have client timeouts, so we treat them as directional
counterexamples rather than stable capacity estimates.

\textbf{Limitations.}
The evidence is limited to Qwen3-8B, one two-NPU platform, one software version,
and a fixed 2,048-token bound. We have not completed best-tuned baseline chunk
sensitivity, controlled same-backend mixed/pure isolation, static-ratio
sensitivity, or cross-model validation. The MaaS input is reconstructed from
aggregates and its tighter SLO is post hoc. The results therefore support an
end-to-end comparison of these evaluated deployments, but do not isolate a
kernel-level cause or establish a universally optimal phase share.
